\section{Introduction}
\label{sec:introduction}

A visitor in a virtual exhibition points at a product panel and asks an embodied
agent, ``What is this?''  The expected interaction appears simple: preserve the
visible selection, contact an agent that is responsible for the selected
object, invoke the grounded-answer operation, and present a response that still
refers to the same object.  Yet those decisions often live in different parts
of an implementation.  Unity stores colliders and scene objects; configuration
files assign knowledge and roles; orchestration code selects an agent; an
endpoint invokes a language model; and logs record only that a request
succeeded.  Each part can be locally valid while the composed interaction is
wrong.

This mismatch is particularly consequential in an intelligent interface.  A
language model can produce a plausible sentence even when the interface has
lost the user's target, routed to an unrelated agent, or invoked an operation
that was never declared for that target.  Transport success therefore does not
establish interaction success.  The interface needs a boundary that answers
four deterministic questions before probabilistic generation begins:
\emph{Which modeled entity was selected?  Which agent is responsible?  Which
capability may be used?  Which concrete action realizes it?}  After generation,
the same boundary must determine whether the returned evidence still agrees
with those decisions.

We investigate an \emph{executable interaction contract} for this purpose.
The contract connects a user-facing scenario and spatial target to a provider,
capability use, runtime binding, runtime action, expected assertion, and
observation.  Its current realization, \systemname{} (Functional
Meta-Language-defined Structure), extends an MLDS-style machine-readable
artifact with functional interaction relations
\cite{fischer2025machinereadable}.  A generation pipeline materializes the
contract for a Unity desktop client and a multi-agent backend.  The client
provides a ray-selected target candidate, but the backend treats the pinned
model as authoritative: it validates the candidate's stable identity, derives
its group and zone, resolves one responsible provider, checks a permitted
one-hop route, and selects one target- and provider-specific runtime binding.
No endpoint name supplied by the client can substitute for that path.

We call a request \emph{contract-grounded} when its target, provider,
capability use, capability, binding, action, and expected assertion resolve
within one pinned model version.  The guarantee is conditional but concrete.
For an admitted request, the implemented preflight found one complete,
model-local path and the response evidence agrees with the trusted target and
provider.  A request-decidable violation stops before a response-model call or
session mutation.  A response-dependent violation, such as a generated handoff
to an undeclared destination, causes transactional rollback.  In both cases the
system can report the enforcement stage and violated relation rather than a
generic chat failure.

This guarantee deliberately begins \emph{after} perceptual target selection.
It does not infer an unobserved intention, prove that ray geometry is honest,
or establish factual correctness of generated language.  A multimodal
reference resolver can propose the candidate; the contract determines whether
that candidate has an admissible path through the interface.  The resulting
separation is useful because reference inference and action permission fail in
different ways and require different evidence.

The work is motivated by three interaction requirements drawn from the
implemented visitor path and established human--AI interaction guidance
\cite{amershi2019guidelines,yang2020reexamining,weisz2024design}:

\begin{description}
  \item[Stable reference.] The object highlighted by the client remains the
  target while the visitor composes the request, and the response is presented
  only if the returned identifiers match that pending selection.
  \item[Visible unresolved states.] No target and ambiguous target are explicit
  states.  They block a deictic request rather than falling through to a fluent
  ungrounded answer.
  \item[Inspectable responsibility.] The chosen provider, route reason,
  capability, binding, and action are retained as structured evidence for
  developer diagnosis.
\end{description}

We do not infer that these mechanisms improve trust, preference, or usability;
that requires human evaluation.  We instead ask whether the system implements
the claimed control boundary, whether migration preserves the behavior already
encoded in the direct artifacts, and what assurance and cost the additional
relations introduce.

Accordingly, we study three research questions:

\begin{description}
  \item[RQ1---Coverage and non-regression.] Can the authored spatial cases be
  migrated to complete target--provider--capability--action paths while
  preserving their existing ownership and routing decisions under a fixed
  policy?
  \item[RQ2---Fail-closed enforcement.] Does the implemented Unity--backend
  path admit exactly the contract-valid local and one-hop requests, preserve
  target and provider evidence, and reject stale, ambiguous, unreachable, or
  contradictory requests without retained side effects?
  \item[RQ3---Added assurance and cost.] Which cross-layer inconsistencies are
  expressible and localizable in the contract representation beyond the
  direct artifacts, and what representation and local execution costs
  accompany that structure?
\end{description}

The paper makes four contributions:

\begin{enumerate}
  \item an interaction-oriented failure model and executable contract that
  binds a spatial selection to one provider, capability, action, and evidence
  path under a pinned model version;
  \item an implemented fail-closed architecture across Unity and a
  conversational-agent backend, including authoritative server-side
  re-resolution, staged admission checks, one-hop handoff constraints, and
  transactional response validation;
  \item a reproducible, API-free evaluation over three spatial environments,
  comprising complete-chain audits, a normalized direct-wiring
  non-regression comparison, 459 executable routing probes, targeted invalid
  requests, and Unity batch smokes; and
  \item a separated analysis of common faults, contract-only cross-layer
  faults, representation edits, and local structural runtime, exposing both
  the added invariant scope and its cost.
\end{enumerate}

The central claim is not that a metamodel by itself makes an interface
intelligent.  Rather, the explicit relations determine what the intelligent
interface is allowed to interpret and present.  This shifts a class of
wrong-target, wrong-provider, and undeclared-action failures out of
probabilistic dialogue behavior and into a deterministic, inspectable boundary.
